\documentclass[letterpaper,12pt]{article}
\usepackage[margin=0in]{geometry}
\usepackage[T1]{fontenc}
\usepackage{lmodern}
\pdfmapfile{+lm.map}
\usepackage{tikz}
\pagestyle{empty}
\setlength{\parindent}{0pt}
\hyphenpenalty=10000
\exhyphenpenalty=10000
\begin{document}

\null
\begin{tikzpicture}[remember picture,overlay,x=1in,y=1in]
\begin{scope}[shift={(current page.south west)}]
\node[anchor=west,inner sep=0pt,font=\fontsize{11.5}{13}\selectfont] at (0.65,10.5) {Week 15 / Nearest-site regions / Grades 4--5};
\node[anchor=north west,inner sep=0pt,text width=7.2in,align=left,font=\fontsize{11.5}{14}\selectfont] at (0.65,10.10) {Keep the dots fixed and in different places. Use straight-line distance to their centers. A dot's region includes every place where it is nearest, including shared ties. The frame shows part of a larger plane.};
\node[anchor=north west,inner sep=0pt,text width=7.2in,align=left,font=\fontsize{14}{18}\selectfont] at (0.65,9.37) {\textbf{Problem 1:} Draw the two nearest-dot regions. Explain why your boundary works for every point on the sheet.};
\begin{scope}[shift={(4.25,5.3)},x=0.966666667in,y=0.966666667in]
\draw[black!55,line width=0.5pt] (-3,-3) rectangle (3,3);
\draw[black,fill=white,line width=0.65pt] (-1.5,-1) circle[radius=0.062];
\fill[black] (-1.5,-1) circle[radius=0.018];
\node[anchor=south west,inner sep=0pt,font=\bfseries\fontsize{15}{17}\selectfont] at (-1.4,-0.9299999999999999) {A};
\draw[black,fill=white,line width=0.65pt] (1.5,1) circle[radius=0.062];
\fill[black] (1.5,1) circle[radius=0.018];
\node[anchor=south west,inner sep=0pt,font=\bfseries\fontsize{15}{17}\selectfont] at (1.6,1.07) {B};
\end{scope}
\node[anchor=west,inner sep=0pt,font=\fontsize{10}{12}\selectfont] at (0.65,0.48) {Bellingham Math Circle / Week 15 / F15-45-v1};
\node[anchor=east,inner sep=0pt,font=\fontsize{10}{12}\selectfont] at (7.85,0.48) {1};
\end{scope}
\end{tikzpicture}
\newpage
\null
\begin{tikzpicture}[remember picture,overlay,x=1in,y=1in]
\begin{scope}[shift={(current page.south west)}]
\node[anchor=west,inner sep=0pt,font=\fontsize{11.5}{13}\selectfont] at (0.65,10.5) {Week 15 / Nearest-site regions / Grades 4--5};
\node[anchor=north west,inner sep=0pt,text width=7.2in,align=left,font=\fontsize{14}{18}\selectfont] at (0.65,9.85) {\textbf{Problem 2:} Draw the three regions, including all the places shared by nearest dots.};
\begin{scope}[shift={(4.25,5.3)},x=0.966666667in,y=0.966666667in]
\draw[black!55,line width=0.5pt] (-3,-3) rectangle (3,3);
\draw[black,fill=white,line width=0.65pt] (-2,0) circle[radius=0.062];
\fill[black] (-2,0) circle[radius=0.018];
\node[anchor=south west,inner sep=0pt,font=\bfseries\fontsize{15}{17}\selectfont] at (-1.9,0.07) {A};
\draw[black,fill=white,line width=0.65pt] (2,0) circle[radius=0.062];
\fill[black] (2,0) circle[radius=0.018];
\node[anchor=south west,inner sep=0pt,font=\bfseries\fontsize{15}{17}\selectfont] at (2.1,0.07) {B};
\draw[black,fill=white,line width=0.65pt] (0,2) circle[radius=0.062];
\fill[black] (0,2) circle[radius=0.018];
\node[anchor=south west,inner sep=0pt,font=\bfseries\fontsize{15}{17}\selectfont] at (0.1,2.07) {C};
\end{scope}
\node[anchor=west,inner sep=0pt,font=\fontsize{10}{12}\selectfont] at (0.65,0.48) {Bellingham Math Circle / Week 15 / F15-45-v1};
\node[anchor=east,inner sep=0pt,font=\fontsize{10}{12}\selectfont] at (7.85,0.48) {2};
\end{scope}
\end{tikzpicture}
\newpage
\null
\begin{tikzpicture}[remember picture,overlay,x=1in,y=1in]
\begin{scope}[shift={(current page.south west)}]
\node[anchor=west,inner sep=0pt,font=\fontsize{11.5}{13}\selectfont] at (0.65,10.5) {Week 15 / Nearest-site regions / Grades 4--5};
\node[anchor=north west,inner sep=0pt,text width=7.2in,align=left,font=\fontsize{14}{18}\selectfont] at (0.65,9.85) {\textbf{Problem 3:} Draw D's whole region. Can a straight segment joining two points of that region ever leave it? Explain why your answer holds for every pair of points.};
\begin{scope}[shift={(4.25,5.3)},x=0.966666667in,y=0.966666667in]
\draw[black!55,line width=0.5pt] (-3,-3) rectangle (3,3);
\draw[black,fill=white,line width=0.65pt] (-2,-1) circle[radius=0.062];
\fill[black] (-2,-1) circle[radius=0.018];
\node[anchor=south west,inner sep=0pt,font=\bfseries\fontsize{15}{17}\selectfont] at (-1.9,-0.9299999999999999) {A};
\draw[black,fill=white,line width=0.65pt] (2,-1) circle[radius=0.062];
\fill[black] (2,-1) circle[radius=0.018];
\node[anchor=south west,inner sep=0pt,font=\bfseries\fontsize{15}{17}\selectfont] at (2.1,-0.9299999999999999) {B};
\draw[black,fill=white,line width=0.65pt] (0,2) circle[radius=0.062];
\fill[black] (0,2) circle[radius=0.018];
\node[anchor=south west,inner sep=0pt,font=\bfseries\fontsize{15}{17}\selectfont] at (0.1,2.07) {C};
\draw[black,fill=white,line width=0.65pt] (0,0) circle[radius=0.062];
\fill[black] (0,0) circle[radius=0.018];
\node[anchor=south west,inner sep=0pt,font=\bfseries\fontsize{15}{17}\selectfont] at (0.1,0.07) {D};
\end{scope}
\node[anchor=west,inner sep=0pt,font=\fontsize{10}{12}\selectfont] at (0.65,0.48) {Bellingham Math Circle / Week 15 / F15-45-v1};
\node[anchor=east,inner sep=0pt,font=\fontsize{10}{12}\selectfont] at (7.85,0.48) {3};
\end{scope}
\end{tikzpicture}
\newpage
\null
\begin{tikzpicture}[remember picture,overlay,x=1in,y=1in]
\begin{scope}[shift={(current page.south west)}]
\node[anchor=west,inner sep=0pt,font=\fontsize{11.5}{13}\selectfont] at (0.65,10.5) {Week 15 / Nearest-site regions / Grades 4--5};
\node[anchor=north west,inner sep=0pt,text width=7.2in,align=left,font=\fontsize{14}{18}\selectfont] at (0.65,9.85) {\textbf{Problem 4:} Draw the four regions. Find every place shared by three or more nearest dots.};
\begin{scope}[shift={(4.25,5.3)},x=0.966666667in,y=0.966666667in]
\draw[black!55,line width=0.5pt] (-3,-3) rectangle (3,3);
\draw[black,fill=white,line width=0.65pt] (-2,2) circle[radius=0.062];
\fill[black] (-2,2) circle[radius=0.018];
\node[anchor=south west,inner sep=0pt,font=\bfseries\fontsize{15}{17}\selectfont] at (-1.9,2.07) {A};
\draw[black,fill=white,line width=0.65pt] (2,2) circle[radius=0.062];
\fill[black] (2,2) circle[radius=0.018];
\node[anchor=south west,inner sep=0pt,font=\bfseries\fontsize{15}{17}\selectfont] at (2.1,2.07) {B};
\draw[black,fill=white,line width=0.65pt] (2,-2) circle[radius=0.062];
\fill[black] (2,-2) circle[radius=0.018];
\node[anchor=south west,inner sep=0pt,font=\bfseries\fontsize{15}{17}\selectfont] at (2.1,-1.93) {C};
\draw[black,fill=white,line width=0.65pt] (-2,-2) circle[radius=0.062];
\fill[black] (-2,-2) circle[radius=0.018];
\node[anchor=south west,inner sep=0pt,font=\bfseries\fontsize{15}{17}\selectfont] at (-1.9,-1.93) {D};
\end{scope}
\node[anchor=west,inner sep=0pt,font=\fontsize{10}{12}\selectfont] at (0.65,0.48) {Bellingham Math Circle / Week 15 / F15-45-v1};
\node[anchor=east,inner sep=0pt,font=\fontsize{10}{12}\selectfont] at (7.85,0.48) {4};
\end{scope}
\end{tikzpicture}
\newpage
\null
\begin{tikzpicture}[remember picture,overlay,x=1in,y=1in]
\begin{scope}[shift={(current page.south west)}]
\node[anchor=west,inner sep=0pt,font=\fontsize{11.5}{13}\selectfont] at (0.65,10.5) {Week 15 / Nearest-site regions / Grades 4--5};
\node[anchor=north west,inner sep=0pt,text width=7.2in,align=left,font=\fontsize{14}{18}\selectfont] at (0.65,9.85) {\textbf{Problem 5:} E is added to the dots in Problem 4. Draw all five regions and mark the old boundaries that disappear. Can adding a new dot ever give an old dot more places? Explain.};
\begin{scope}[shift={(4.25,5.3)},x=0.966666667in,y=0.966666667in]
\draw[black!55,line width=0.5pt] (-3,-3) rectangle (3,3);
\draw[black,fill=white,line width=0.65pt] (-2,2) circle[radius=0.062];
\fill[black] (-2,2) circle[radius=0.018];
\node[anchor=south west,inner sep=0pt,font=\bfseries\fontsize{15}{17}\selectfont] at (-1.9,2.07) {A};
\draw[black,fill=white,line width=0.65pt] (2,2) circle[radius=0.062];
\fill[black] (2,2) circle[radius=0.018];
\node[anchor=south west,inner sep=0pt,font=\bfseries\fontsize{15}{17}\selectfont] at (2.1,2.07) {B};
\draw[black,fill=white,line width=0.65pt] (2,-2) circle[radius=0.062];
\fill[black] (2,-2) circle[radius=0.018];
\node[anchor=south west,inner sep=0pt,font=\bfseries\fontsize{15}{17}\selectfont] at (2.1,-1.93) {C};
\draw[black,fill=white,line width=0.65pt] (-2,-2) circle[radius=0.062];
\fill[black] (-2,-2) circle[radius=0.018];
\node[anchor=south west,inner sep=0pt,font=\bfseries\fontsize{15}{17}\selectfont] at (-1.9,-1.93) {D};
\draw[black,fill=white,line width=0.65pt] (0,0) circle[radius=0.062];
\fill[black] (0,0) circle[radius=0.018];
\node[anchor=south west,inner sep=0pt,font=\bfseries\fontsize{15}{17}\selectfont] at (0.1,0.07) {E};
\end{scope}
\node[anchor=west,inner sep=0pt,font=\fontsize{10}{12}\selectfont] at (0.65,0.48) {Bellingham Math Circle / Week 15 / F15-45-v1};
\node[anchor=east,inner sep=0pt,font=\fontsize{10}{12}\selectfont] at (7.85,0.48) {5};
\end{scope}
\end{tikzpicture}
\newpage
\null
\begin{tikzpicture}[remember picture,overlay,x=1in,y=1in]
\begin{scope}[shift={(current page.south west)}]
\node[anchor=west,inner sep=0pt,font=\fontsize{11.5}{13}\selectfont] at (0.65,10.5) {Week 15 / Nearest-site regions / Grades 4--5};
\node[anchor=north west,inner sep=0pt,text width=7.2in,align=left,font=\fontsize{14}{18}\selectfont] at (0.65,9.85) {\textbf{Problem 6:} Place as few new dots as you can so B's entire region fits inside the frame. Draw that region and explain why fewer new dots cannot work.};
\begin{scope}[shift={(4.25,5.3)},x=0.966666667in,y=0.966666667in]
\draw[black!55,line width=0.5pt] (-3,-3) rectangle (3,3);
\draw[black,fill=white,line width=0.65pt] (-2,0) circle[radius=0.062];
\fill[black] (-2,0) circle[radius=0.018];
\node[anchor=south west,inner sep=0pt,font=\bfseries\fontsize{15}{17}\selectfont] at (-1.9,0.07) {A};
\draw[black,fill=white,line width=0.65pt] (0,0) circle[radius=0.062];
\fill[black] (0,0) circle[radius=0.018];
\node[anchor=south west,inner sep=0pt,font=\bfseries\fontsize{15}{17}\selectfont] at (0.1,0.07) {B};
\draw[black,fill=white,line width=0.65pt] (2,0) circle[radius=0.062];
\fill[black] (2,0) circle[radius=0.018];
\node[anchor=south west,inner sep=0pt,font=\bfseries\fontsize{15}{17}\selectfont] at (2.1,0.07) {C};
\end{scope}
\node[anchor=west,inner sep=0pt,font=\fontsize{10}{12}\selectfont] at (0.65,0.48) {Bellingham Math Circle / Week 15 / F15-45-v1};
\node[anchor=east,inner sep=0pt,font=\fontsize{10}{12}\selectfont] at (7.85,0.48) {6};
\end{scope}
\end{tikzpicture}
\newpage
\null
\begin{tikzpicture}[remember picture,overlay,x=1in,y=1in]
\begin{scope}[shift={(current page.south west)}]
\node[anchor=west,inner sep=0pt,font=\fontsize{11.5}{13}\selectfont] at (0.65,10.5) {Week 15 / Nearest-site regions / Grades 4--5};
\node[anchor=north west,inner sep=0pt,text width=7.2in,align=left,font=\fontsize{14}{18}\selectfont] at (0.65,9.85) {\textbf{Problem 7:} Place three new dots so A's whole region is exactly the shaded triangle.};
\begin{scope}[shift={(4.25,5.3)},x=0.966666667in,y=0.966666667in]
\fill[black!8] (-2,-1)--(2,-1)--(0,2)--cycle;
\draw[black,line width=0.6pt,dash pattern=on 4pt off 3pt] (-2,-1)--(2,-1)--(0,2)--cycle;
\draw[black!55,line width=0.5pt] (-3,-3) rectangle (3,3);
\draw[black,fill=white,line width=0.65pt] (0,0) circle[radius=0.062];
\fill[black] (0,0) circle[radius=0.018];
\node[anchor=south west,inner sep=0pt,font=\bfseries\fontsize{15}{17}\selectfont] at (0.1,0.07) {A};
\end{scope}
\node[anchor=west,inner sep=0pt,font=\fontsize{10}{12}\selectfont] at (0.65,0.48) {Bellingham Math Circle / Week 15 / F15-45-v1};
\node[anchor=east,inner sep=0pt,font=\fontsize{10}{12}\selectfont] at (7.85,0.48) {7};
\end{scope}
\end{tikzpicture}
\newpage
\null
\begin{tikzpicture}[remember picture,overlay,x=1in,y=1in]
\begin{scope}[shift={(current page.south west)}]
\node[anchor=west,inner sep=0pt,font=\fontsize{11.5}{13}\selectfont] at (0.65,10.5) {Week 15 / Nearest-site regions / Grades 4--5};
\node[anchor=north west,inner sep=0pt,text width=7.2in,align=left,font=\fontsize{14}{18}\selectfont] at (0.65,9.85) {\textbf{Problem 8:} For R and S separately, decide whether new dots can take that point away from A while P and Q still belong to A. Draw a placement when it is possible; explain when it is impossible.};
\begin{scope}[shift={(4.25,5.3)},x=0.966666667in,y=0.966666667in]
\draw[black!55,line width=0.5pt] (-3,-3) rectangle (3,3);
\draw[black,line width=0.65pt] (-2.045,0)--(-1.955,0) (-2,-0.045)--(-2,0.045);
\node[anchor=west,inner sep=0pt,font=\fontsize{13}{15}\selectfont] at (-1.89,0.02) {P};
\draw[black,line width=0.65pt] (1.955,0)--(2.045,0) (2,-0.045)--(2,0.045);
\node[anchor=west,inner sep=0pt,font=\fontsize{13}{15}\selectfont] at (2.11,0.02) {Q};
\draw[black,line width=0.65pt] (-0.045,0)--(0.045,0) (0,-0.045)--(0,0.045);
\node[anchor=west,inner sep=0pt,font=\fontsize{13}{15}\selectfont] at (0.11,0.02) {R};
\draw[black,line width=0.65pt] (-0.045,-2)--(0.045,-2) (0,-2.045)--(0,-1.955);
\node[anchor=west,inner sep=0pt,font=\fontsize{13}{15}\selectfont] at (0.11,-1.98) {S};
\draw[black,fill=white,line width=0.65pt] (0,2) circle[radius=0.062];
\fill[black] (0,2) circle[radius=0.018];
\node[anchor=south west,inner sep=0pt,font=\bfseries\fontsize{15}{17}\selectfont] at (0.1,2.07) {A};
\end{scope}
\node[anchor=west,inner sep=0pt,font=\fontsize{10}{12}\selectfont] at (0.65,0.48) {Bellingham Math Circle / Week 15 / F15-45-v1};
\node[anchor=east,inner sep=0pt,font=\fontsize{10}{12}\selectfont] at (7.85,0.48) {8};
\end{scope}
\end{tikzpicture}
\end{document}